% ==============================================================================
% CAPÍTULO 2 - METODOLOGIA
% ==============================================================================

\section{Metodologia}
\label{sec:metodologia}

Esta seção descreve o delineamento da pesquisa: sua classificação, os
procedimentos seguidos, as métricas e os testes estatísticos adotados, a
mudança de planejamento ocorrida durante o experimento e o papel das
ferramentas de inteligência artificial no método. A execução técnica de cada
etapa é detalhada na \autoref{sec:desenvolvimento}.

\subsection{Classificação da pesquisa}
\label{sec:classificacao}

Quanto à natureza, trata-se de pesquisa aplicada, pois busca responder a uma
pergunta prática da engenharia de software com LLMs: se e como vale a pena
investir em um \textit{harness} de repositório. Quanto à abordagem, é
quantitativa: todas as variáveis de resposta são medidas numericamente e
analisadas com estatística descritiva e inferencial. Quanto aos objetivos, é
explicativa, no sentido de \citeonline{gil2017}, pois procura identificar o
efeito de um fator manipulado sobre variáveis de resposta. Quanto aos
procedimentos, é um experimento controlado, na terminologia de
\citeonline{wohlin2012}, organizado como um estudo de ablação: a partir de um
\textit{harness} completo, camadas são removidas progressivamente e o efeito de
cada remoção é medido.

O fator manipulado é o nível do \textit{harness}, com quatro valores
cumulativos (\texttt{none}, \texttt{agents}, \texttt{skills} e \texttt{full},
descritos na \autoref{sec:objeto}). O modelo de linguagem é um fator de bloco:
as comparações entre níveis são feitas dentro de cada modelo, nunca entre
modelos, porque os modelos foram executados por \textit{runners} diferentes. A
tarefa é a unidade de pareamento: cada tarefa é executada por cada modelo em
cada nível, de modo que o efeito de um nível é estimado pela diferença, tarefa a
tarefa, em relação ao nível \texttt{none}. Uma célula do experimento é,
portanto, a tripla modelo~$\times$~tarefa~$\times$~nível, e uma execução
(\textit{run}) é a materialização de uma célula.

\subsection{Procedimentos}
\label{sec:procedimentos}

O experimento seguiu oito etapas, executadas entre 26 e 29 de outubro de 2026.

\textbf{Seleção das tarefas a partir de PRs.} O repositório-sujeito é uma
aplicação web de gestão de filas de atendimento escrita em Laravel (PHP) com
Bulma CSS e Livewire, privada e desenvolvida pelo autor. Cada tarefa reproduz
um \textit{pull request} (PR) ou \textit{branch} desse repositório: a árvore
no \textit{commit} base do PR é o ponto de partida do agente, um enunciado
escrito a partir do PR e da documentação de domínio é a tarefa, e os arquivos
de teste adicionados ou modificados pelo PR são o oráculo oculto, copiados para
a árvore apenas no momento da avaliação. O enunciado carrega o contrato público
de que os testes dependem (nomes de classes, métodos, rotas, chaves de
configuração), para que a tarefa meça engenharia e não adivinhação de nomes;
os pontos em que o acoplamento permanece forte são registrados por tarefa.
Dois PRs grandes, sem \textit{commit} intermediário aproveitável, foram
mantidos como tarefas compostas de referência e divididos em subtarefas que
compartilham base e solução, mas avaliam subconjuntos separáveis dos testes.
O conjunto final tem 14 tarefas: nove PRs ou \textit{branches} inteiros e
cinco subtarefas das duas compostas (\autoref{sec:benchmark} e Apêndice~A).

\textbf{Congelamento do \textit{harness}.} O \textit{harness} estudado foi
extraído de uma \textit{tag} anotada do repositório-sujeito
(\texttt{tcc-baseline-v1}, de 26 de outubro de 2026) e copiado para o
repositório do \textit{benchmark} em 27 de outubro, byte a byte. Todas as
tarefas, independentemente do ponto da história de onde vêm, recebem esse mesmo
\textit{snapshot}, o que elimina a variação do \textit{harness} entre tarefas e
isola o fator estudado.

\textbf{Definição dos níveis.} Os quatro níveis são cumulativos: \texttt{none}
não contém nenhum artefato do \textit{harness}; \texttt{agents} adiciona as
regras sempre ativas (\texttt{AGENTS.md} e \texttt{CLAUDE.md}); \texttt{skills}
adiciona o diretório \texttt{ai/} com \textit{skills}, comandos, regras e
modelos de código para leitura sob demanda; \texttt{full} adiciona a
configuração dos \textit{runners} e dos servidores de linguagem
(\texttt{opencode.json} e \texttt{.claude/settings.json}). O \textit{harness} é
aplicado em um \textit{commit} separado, marcado com a \textit{tag}
\texttt{bench/start}, de modo que o \textit{diff} do agente nunca inclui os
arquivos do \textit{harness}.

\textbf{Modelos e \textit{runners}.} Foram avaliados seis modelos de escalas e
origens distintas. Claude Sonnet 5.5 e Claude Haiku 5.5 foram executados
exclusivamente pelo Claude Code 2.1.312, com esforço de raciocínio
\texttt{high}. Muse Spark 1.3, DeepSeek v4.1 flash e Beam (plano gratuito) foram
executados pelo OpenCode v2.2.0 por meio do provedor OpenCode Go, com os
parâmetros padrão do provedor. O gpt-oss-20b \cite{openai2025gptoss} foi
executado localmente, pelo OpenCode contra um servidor llama.cpp
\cite{llamacpp2025} em uma GPU de 16~GB, com quantização MXFP4, contexto de
128 mil \textit{tokens} e a amostragem recomendada pelo fabricante. Os modelos
da Anthropic rodaram apenas no Claude Code, nunca por intermédio do OpenCode,
para que o \textit{runner} nativo de cada família fosse usado. Em cada
execução, as mensagens da sessão foram conferidas para garantir que somente o
modelo sob teste respondeu.

\textbf{Execução orquestrada e avaliação automática.} Um orquestrador escrito
em Python para este trabalho prepara cada célula, executa o agente em um
contêiner Docker com acesso apenas à própria árvore e avalia o resultado em um
segundo contêiner, copiando os testes ocultos e executando o Pint, a suíte
completa e os testes ocultos isoladamente. Todos os agentes recebem o mesmo
\textit{prompt}: um preâmbulo não interativo fixo seguido do enunciado sem
alterações, com 30 minutos de limite; a ordem das células foi embaralhada com
semente fixa (20261028) e cada execução registra versões, imagens e o
\textit{hash} do \textit{prompt}. A \autoref{sec:orquestrador} detalha o
orquestrador.

\textbf{Análise de qualidade em três camadas.} Como a taxa de acerto ficou
próxima do teto para a maioria dos modelos, uma segunda fase perguntou se o
\textit{harness} muda \emph{como} o código é escrito. Essa fase trabalha sobre
os \textit{patches} e as árvores deixadas pelas execuções, sem rodar agentes de
novo, em três camadas: (1) métricas estáticas, com Larastan nível 5
\cite{larastan2025} para contar erros novos em relação à base, complexidade
ciclomática de McCabe por função, casos de teste escritos e distância da
solução de referência; (2) dez regras heurísticas (C1 a C10) derivadas do
próprio \texttt{AGENTS.md}, cada uma aplicada apenas quando a mudança toca o
seu assunto; (3) revisão às cegas por até três juízes LLM por mudança, Claude
Opus 5.5, Muse Spark 1.3 e DeepSeek v4.1 flash, seguindo uma rubrica de sete
dimensões com notas de 1 a 5 (Apêndice~B). Cada mudança recebeu dois juízes de
famílias distintas da do modelo que a escreveu e, quando existe, o juiz da
própria família; a nota da execução é a média de todos os seus juízes, e a
nota do juiz da família é comparada à dos dois de fora como controle de
autopreferência. O \autoref{qdr:camadas} resume as três camadas. Um piloto com
duas tarefas (46 execuções, 92 revisões) validou o método antes da rodada
completa (309 execuções, 836 revisões), e todas as
violações da camada 2 apontadas no piloto foram lidas manualmente contra o
código.

\begin{quadro}[htb]
    \centering
    \ABNTEXfontereduzida
    \caption{\label{qdr:camadas}As três camadas da análise de qualidade de código}
    \begin{tabular}{|l|p{4.6cm}|p{3.4cm}|p{4.2cm}|}
        \hline
        \textbf{Camada} & \textbf{O que mede} & \textbf{Ferramenta} & \textbf{Saída por execução} \\
        \hline
        1. Estática & Erros novos de análise estática em relação à base; complexidade ciclomática por função; casos de teste e linhas escritas, separadas em produção e testes; distância da solução de referência & Larastan nível 5, analisador de complexidade em PHP, contagem de testes Pest, índice de Jaccard dos arquivos de produção & \texttt{larastan\_new\_errors}, \texttt{cc\_mean}, \texttt{test\_cases\_added}, \texttt{ref\_files\_jaccard} \\
        \hline
        2. Convenções & Dez regras do \texttt{AGENTS.md} (C1 a C10), cada uma aplicada só quando a mudança toca o seu assunto & Heurísticas sobre o \textit{patch}, com evidência guardada por violação & Escore de 0 a 1 (fração das regras aplicáveis cumpridas) \\
        \hline
        3. Juízes & Sete dimensões (requisitos, convenções, segurança, testes, simplicidade, legibilidade, escopo) e nota geral, de 1 a 5, com evidência obrigatória & Até três juízes LLM às cegas (Opus 5.5, Muse Spark 1.3, DeepSeek v4.1 flash): dois de famílias distintas da do autor e, quando existe, o da própria família; rubrica e esquema JSON & \texttt{review.json} por juiz; média de todos os juízes por dimensão \\
        \hline
    \end{tabular}
    \fonteautor
\end{quadro}

\textbf{Análise estatística.} Descrita na \autoref{sec:estatistica}.

\subsection{Métricas}
\label{sec:metricas}

A métrica principal de acerto é o \textit{pass@1}: a célula passa quando todos
os testes ocultos passam em uma única tentativa, na definição de
\citeonline{chen2021codex} para $k = 1$. Como métrica contínua, registra-se a
fração de testes ocultos aprovados. Regressões são testes fora do conjunto
oculto que falham além dos que já falhavam na árvore intocada, medidos por
tarefa com um agente nulo como linha de base. O custo é medido por turnos do
agente, \textit{tokens} de entrada (incluindo leitura de \textit{cache}) e de
saída e tempo de parede; turnos e \textit{tokens} são comparáveis entre níveis
de um mesmo modelo, enquanto o tempo depende do provedor e da concorrência. A
qualidade é medida pelas três camadas descritas acima: escore de convenções
(fração das regras aplicáveis cumpridas), erros novos do Larastan,
complexidade média, casos de teste novos, linhas inseridas em código de
produção e em testes, índice de Jaccard entre os arquivos
de produção tocados e os da referência, e as notas dos juízes em cada dimensão
e na nota geral, com a média de todos os juízes da execução.

\subsection{Análise estatística}
\label{sec:estatistica}

Todas as comparações são pareadas por tarefa, dentro de cada modelo, contra o
nível \texttt{none}. Para o \textit{pass@1}, um resultado binário, usa-se o teste
exato de McNemar \cite{mcnemar1947}, que considera apenas as tarefas
discordantes entre os dois níveis; os intervalos de confiança de 95\% das
proporções são de Wilson \cite{wilson1927}. Para as métricas contínuas e ordinais
de qualidade, usa-se o teste exato dos postos sinalizados de Wilcoxon
\cite{wilcoxon1945}, adequado a amostras pequenas sem suposição de normalidade.
A concordância entre juízes é medida pelo $\kappa$ ponderado de
\citeonline{cohen1968}, com pesos quadráticos, interpretado pela escala de
\citeonline{landis1977}, e pela correlação de postos de Spearman
\cite{spearman1904}. A autopreferência é medida, em cada mudança, pela
diferença entre a nota do juiz da família do autor e a média dos dois juízes de
fora (Wilcoxon exato pareado), descontada a leniência geral desse juiz nas
mudanças de outras famílias. A validade das notas dos juízes é verificada pela
correlação de Spearman entre a nota e medidas independentes (acerto nos testes
ocultos, escore de convenções e erros do Larastan). O agregado entre modelos
usa os 70 pares por nível dos cinco modelos com a matriz completa (5 modelos
$\times$ 14 tarefas), excluindo o Beam, que ficou incompleto. Não foi aplicada
correção para comparações múltiplas; o número de comparações é informado e a
leitura dos $p$-valores considera esse fato (\autoref{sec:ameacas}).

\subsection{Mudança de planejamento: da rodada completa à rodada reduzida}
\label{sec:mudanca}

O plano original previa cinco modelos, 14 tarefas, quatro níveis e três
repetições (840 células), com os modelos GLM 5.3 flash e MiMo v2.6 flash no
lugar de Muse Spark, DeepSeek e Beam. A rodada foi iniciada em 28 de outubro
de 2026 às 18h51 e cancelada às 19h40, depois de 41 células, quando a
estimativa de duração com tempos reais ultrapassou 30 horas; o gargalo eram os
dois modelos mais lentos do provedor, com cerca de 75 e 45 \textit{tokens} por
segundo em sessões agênticas. Com o limite de 12 horas disponível, optou-se
por uma rodada reduzida: uma repetição em vez de três, os três modelos mais
rápidos do provedor (Muse Spark, Beam e DeepSeek, medidos em 376, 339 e 280
\textit{tokens} por segundo isoladamente) e os dois modelos locais. Essa
decisão foi registrada com data e justificativa no diário do orquestrador. A
consequência metodológica é que a repetição vem das tarefas, não de execuções
repetidas da mesma célula: a variação entre tarefas tende a superar a variação
entre execuções da mesma tarefa, mas a variância de um mesmo modelo na mesma
célula não é estimada. Durante a rodada reduzida, o segundo modelo local (Qwen
3.6 35B-A3B) foi retirado por lentidão (cerca de 23 minutos por célula, o que
levaria 20 horas) e o provedor do Beam ficou fora do ar; o Beam foi encerrado
com 29 das 56 células e é apresentado apenas de forma descritiva. O desenho
final tem 6 modelos $\times$ 14 tarefas $\times$ 4 níveis $\times$ 1 repetição
(336 células).

\subsection{Uso de inteligência artificial no método}
\label{sec:uso-ia}

A inteligência artificial participa deste trabalho de três maneiras, que
convém distinguir. Primeiro, os modelos de linguagem são o \emph{objeto} do
experimento: os seis modelos avaliados escrevem o código que é medido.
Segundo, modelos são parte do \emph{instrumento} de medida: Claude Opus 5.5,
Muse Spark 1.3 e DeepSeek v4.1 flash atuam como juízes na terceira camada da
análise de qualidade, sob as salvaguardas descritas na
\autoref{sec:qualidade-metodo}. Terceiro, modelos foram \emph{ferramentas} de
construção: o orquestrador, os \textit{scripts} de análise e os relatórios
intermediários foram desenvolvidos com o auxílio do Claude Code (Claude Opus
5.5), sob direção e revisão do autor; o código do repositório-sujeito foi
escrito, em parte, com o Muse Spark, o que é discutido como ameaça à validade;
e o texto deste artigo foi redigido com o auxílio do Claude Fable 5.1 e revisado
pelo autor, que responde integralmente pelo conteúdo. A declaração formal de uso
consta ao fim do trabalho.
